\chapter{Конечный Пейдж--Вуттерс--Стайнспринг-мост}
\label{ch:v8-page-wootters-stinespring-history-gate}

\section{От канала к истории}

Минимальный канал межсекторных стрелок при контрольном значении $p=1/12$ имеет
тринадцать операторов Крауса $K_a$ и точную полноту
\begin{equation}
 \sum_{a=0}^{12}K_a^*K_a=I_{21}.
 \label{eq:v8-pw-kraus-completeness}
\end{equation}
Поэтому он задаёт изометрию Стайнспринга
\begin{equation}
 W\psi=\sum_{a=0}^{12}K_a\psi\otimes|a\rangle,
 \qquad W^*W=I_{21}.
 \label{eq:v8-pw-stinespring-isometry}
\end{equation}

Для двух последовательных свежих сред введём часы с тремя показаниями
$|0\rangle,|1\rangle,|2\rangle$. После вложения всех временных срезов в
общий дополненный носитель определим
\begin{equation}
 |\mathsf H_2(\psi)\rangle=
 \frac1{\sqrt3}\sum_{n=0}^{2}|n\rangle_C\otimes W_n\cdots W_1\psi.
 \label{eq:v8-pw-history-state}
\end{equation}
Это не внешняя траектория, а один коррелированный вектор, содержащий три
условных среза.

\section{Точное условное восстановление}

LCF-ядро разворачивает все слова Крауса. Их верхние числа равны
\begin{equation}
 1,\qquad13,\qquad13^2=169.
\end{equation}
Для начального чистого концевого состояния точное суммирование ветвей даёт
\begin{equation}
 \Tr_{E_1\cdots E_n}
 \left({}_C\langle n|\mathsf H_2\rangle
       \langle\mathsf H_2|n\rangle_C\right)
 =\frac13\Phi_{p,*}^{,n}(\rho_0),
 \qquad n=0,1,2.
 \label{eq:v8-pw-conditional-recovery}
\end{equation}
После условной нормировки каждый срез имеет след один. Все матричные остатки
равны нулю точно, без численного допуска.

Дополненный носитель данных имеет размерность
\begin{equation}
 21\cdot13^2=3549,
\end{equation}
а полный носитель «часы--данные» --- $3\cdot3549=10647$. Рост не является
новой частицей: это явный учёт двух записей среды.

\section{Стационарный родитель и его граница}

Из ограничений распространения
\begin{equation}
 \psi_{n+1}=W_{n+1}\psi_n
\end{equation}
строится положительный без фрустрации родитель истории
$H_{\rm hist}=\sum_nB_n^*B_n$. Для любого исходного вектора
\begin{equation}
 H_{\rm hist}|\mathsf H_2(\psi)\rangle=0.
 \label{eq:v8-pw-stationary-history}
\end{equation}
Тем самым получено $21$-мерное семейство глобально стационарных историй,
условные срезы которых дают дискретную открытую динамику.

Однако это ещё не каноническое аддитивное ограничение
$H_C+H_S=0$. Изометрия $W:\mathbb C^{21}\to\mathbb C^{273}$ фиксирует
только $21$-мерный образ. Для её продолжения до полного унитарного такта
остаётся дополнение размерности
\begin{equation}
 273-21=252,
\end{equation}
на котором допустимо семейство $U(252)$, имеющее $252^2=63504$
вещественных параметра. Канал не выбирает один такой унитарий.

\section{Связь с непрерывным пределом}

Условная история совместима с уже доказанным пределом
\begin{equation}
 \Phi_{u/n,*}^{,n}\longrightarrow e^{u\mathcal L_*}
\end{equation}
в операторной норме. Таким образом, одна цепь теперь связывает стационарную
историю, дискретные показания часов и непрерывную безразмерную полугруппу.
Но источник свежих сред, канонический полный унитарий и физический перевод
$u\mapsto t$ остаются открытыми.

\begin{theorem}[Конечный условный мост истории]
Минимальная Стайнспринг-дилатация Тома VIII допускает точную конечную
стационарную конструкцию истории. Условие на показание часов $n=0,1,2$ и
след по среде восстанавливают $\Phi_{p,*}^{,n}$ без остатка. В слабом
пределе эти срезы согласованы с полугруппой $e^{u\mathcal L_*}$.
\end{theorem}

\begin{nogocriterion}[Канал ещё не является автономными часами]
Нельзя объявлять построенный родитель истории полным механизмом
Пейдж--Вуттерс или физическим временем. Канал не выбирает единственное
унитарное продолжение на $252$-мерном дополнении, не выводит аддитивное
ограничение $H_C+H_S=0$, источник свежих сред и размерную длительность
такта.
\end{nogocriterion}

\begin{protocol}
\begin{enumerate}
 \item размерность системы: \textbf{$21$};
 \item размерность среды одного такта: \textbf{$13$};
 \item показания часов: \textbf{$0,1,2$};
 \item числа ветвей: \textbf{$1,13,169$};
 \item условное восстановление $\Phi_*^n$: \textbf{точно};
 \item след каждого нормированного среза: \textbf{$1$};
 \item стационарный родитель истории: \textbf{получен};
 \item размерность семейства нулевых историй: \textbf{$21$};
 \item полный носитель «часы--данные»: \textbf{$10647$};
 \item свободное дополнение унитария: \textbf{$252$};
 \item канонический полный унитарный такт: \textbf{не получен};
 \item аддитивное ограничение Пейджа--Вуттерса: \textbf{не получено};
 \item предел слабых столкновений: \textbf{унаследован точно};
 \item физическая секунда и стрела: \textbf{не выведены};
 \item следующий тест: \textbf{происхождение или запрет канонического
 автономного унитарного оператора часов}.
\end{enumerate}
\end{protocol}

Машинный сертификат сохранён в
\path{s2t_v8_page_wootters_stinespring_history_gate_results.json}.

\begin{thebibliography}{9}
\bibitem{v8-pw-original}
D.~N.~Page, W.~K.~Wootters,
\emph{Evolution without evolution: Dynamics described by stationary
observables}, Physical Review D 27 (1983), 2885--2892.
\bibitem{v8-pw-attal}
S.~Attal, Y.~Pautrat,
\emph{From repeated to continuous quantum interactions},
Annales Henri Poincar\'e 7 (2006), 59--104.
\end{thebibliography}